\section{Introduction}
An Agent receives a support ticket, checks an order and submits a refund through a third-party commerce service. An egress monitor sees which services are contacted, encrypted message sizes and the gaps between requests. A defender who wants to hide the refund adds decoy calls. Classical transport cover would discard those decoys at a cooperating tunnel endpoint~\cite{apthorpe2019stp}, but the commerce service is not under the defender's control and offers no such endpoint. A relay could move the discard point, yet the hop from the relay to the service stays visible to whoever watches it, and the service itself sees every request. Against such observers the decoys must arrive as real requests and execute. Each one now produces two outputs: an observable communication trace and work performed by the application, which may consume quota, touch other objects or enter the planner's context. Evaluating the trace alone leaves the second output unmeasured.

\begin{figure*}[t]
\centering
\begin{tikzpicture}[x=1cm,y=1cm,font=\small,>=Latex]
\tikzset{box/.style={draw,rounded corners,align=center,minimum height=.85cm,text width=2.4cm},
 note/.style={align=center,font=\footnotesize,text width=3.1cm}}
\node[box,fill=blue!5] (u) at (1.3,1.15) {User\\``Process the next ticket''};
\node[box,fill=blue!5] (a) at (5.1,1.15) {Agent\\plan, check, act, retry};
\node[box,fill=green!6] (m) at (9.2,2.6) {Model service\\next action / streamed output};
\node[box,fill=green!6] (t) at (9.2,-.35) {Tool service\\ticket, order, refund};
\node[box,fill=orange!8] (d) at (14.1,1.15) {Specialist Agent\\optional delegation};
\draw[<->,thick] (u)--node[above,yshift=15pt]{client link C}(a);
\draw[<->,thick] (a.north east)--node[above,sloped]{model link M}(m.west);
\draw[<->,thick] (a.south east)--node[below,sloped]{tool link T}(t.west);
\draw[<->,thick] (a.east)--node[above]{coordination A}(d.west);
\node[note] at (2.0,-.35) {\textbf{Transit observer}\\lengths, gaps, connections;\\destinations when available};
\node[note] at (13.7,-.35) {\textbf{Recipient / active service}\\received arguments;\\optionally selected responses};
\node[draw=black!45,fill=black!3,align=left,text width=15.7cm,inner sep=6pt] at (7.9,-1.65) {
\textbf{Joint protection decision:} Which observed traces remain distinguishable? Does cover preserve the requested goal, allowed business state and resource budget? State where the cover stops and whether its feedback reaches the planner.};
\end{tikzpicture}
\caption{Running support-ticket example and observation rights. A tool result returns to the Agent and can enter the next model request. Each encrypted link exposes its own traffic shape; a recipient additionally sees the request addressed to it. Concurrent tickets may share a connection on either link. The arrows represent application communication, not a one-to-one mapping from tool events to packets.}
\label{fig:chain}
\end{figure*}


\paragraph{The research question}
How should Agent communication defenses be evaluated when cover is executable work? Agent side-channel research is young: our bounded retrieval finds fifteen direct studies, most of them preprints. This early SoK therefore does not estimate field-wide frequencies. It fixes the evaluation object before practice settles: it classifies cover by its termination point, state access, feedback visibility and dispatcher, states when cover must execute (Section~\ref{sec:taxonomy}), maps all forty-four existing settings onto that object, and tests the resulting contract on live executions.

\textbf{A privacy result needs outcomes from the same protected execution.} WRA's published research decoys reduce query recovery from .770 to .565, while its report-utility scores 8.55 and 8.54 belong to a different intervention that changes destination visibility~\cite{jeong2026wra}. The communication-graph study refits an attacker on protected OHTTP captures and reports traffic cost, but not original-task utility under that intervention~\cite{dangol2026workflow}. In the four query/workflow privacy settings that exist so far, none reports both.

\textbf{The process generating cover is itself observable.} Honest tool calls follow model decisions; a naive controller can insert cover without that wait. We use such immediate cover as a calibration baseline: given call boundaries, label-free gap removal restores recognition to 112/120 and 108/120. The substantive tests concern history-matched dispatch, which removes the shortcut but still leaves 93--106/120 recognition under adapted attacks, and a complete shadow planner, whose strongest tested attack sits at the 7/12 pair-ambiguity floor that one independent shadow permits. Task, state and costs are measured on the same runs.

WF and ML-security work already requires realistic observation units, protected-condition attacks and matched costs~\cite{juarez2014critical,mathews2023critical,arp2022dos,tramer2020adaptive}. Agent execution adds three requirements: say where cover terminates, what it may change and whether its results return to the planner.

\paragraph{Contributions}
\begin{itemize}
\item \textbf{A cover taxonomy and necessity rule} (Section~\ref{sec:taxonomy}): cover policies classified by termination, state, feedback and dispatcher, with the deployment conditions under which cover must execute.
\item \textbf{An execution-aware field map and joint contract} (Sections~\ref{sec:correspondence}--\ref{sec:defenses}, \ref{sec:reporting}): all forty-four settings, their transfer relations, and a checklist that binds residual inference, task completion, business invariants and resource cost to one intervention.
\item \textbf{A live test of the contract} (Section~\ref{sec:experiment}): controller dispatch is removable, history matching is not sufficient, and the strongest tested attack on planner cover is explained by a pair-ambiguity floor; we derive how that floor and the model work scale with the number of shadows.
\end{itemize}
